\lecture{14}
\title{Convolutional Neural Networks}

\begin{document}

\subsection{Motivating Convolutional Neural Networks}

\begin{remark}
    These notes assume familiarity with \emph{convolutions},
    specifically for 2D images.  See \href{/6.300/}{6.300} for reference.
\end{remark}

Suppose we want to build a neural network
for classifying images.
How many parameters would we need?

\begin{problem}
    Suppose we take a $1000\,\mathrm{px} \times 1000\,\mathrm{px}$ image
    as the input to a neural network.
    If the first hidden layer has $1000$ features,
    then how many parameters connect
    the input layer to the first hidden layer?
\end{problem}

\begin{solution}
    The input layer has $10^6$ nodes, and the hidden layer has $10^3$ nodes,
    so the neural network between the input layer and the first hidden layer
    appears like a $K_{10^6, 10^3}$ graph.
    That's a total of $\boxed{10^9}$ parameters, one for each edge.

    (And that's assuming the image is grayscale;
    adding color complicates things even more!)
\end{solution}

That's too many to feasibly perform backpropagation on.
The issue isn't the number of nodes in the layers---it's
the number of parameters needed to describe
the relationship between them.

Maybe we don't need the full $K_{10^6, 10^3}$ graph, though.
Or maybe we don't need a unique parameter for every edge
connecting an input node to
the pre-activation of a node in the next layer.

Maybe there's a meaningful subfamily of
linear maps on 2D images
that requires far fewer parameters to specify\dots

\subsection{Convolutions and Pooling}

Technically, these 2D images aren't \emph{really} 2D.
Most images have color channels, too---three of them (RGB), to be precise.
So more generally, we should think of each layer of our NN
as being a three-dimensional $H \times W \times C$ grid of nodes.

The input layer might be a $1000 \times 1000 \times 3$ grid, for example.
Importantly, our hidden layers might also have more than one channel.
So our maps between layers must be able to
meaningfully accept three-dimensional \emph{tensors} as input.

We now introduce our two maps
from 3D tensors to 3D tensors:
pooling and convolving.

\textbf{Definition.}
Given an $H \times W \times C$ image,
we apply \textbf{pooling} to downsample the image
down to $H' \times W' \times C$.
\begin{center}
    \frame{\includegraphics[width=0.4\textwidth]{lecture-14/max-pooling}}
\end{center}
The image above depicts \textbf{max pooling}.
There's also \textbf{average pooling}---you
can probably imagine how that looks.

\begin{remark}
    Pooling always treats each channel independently of every other.
    It may be more accurate to view pooling
    as a map from 2D tensors to smaller 2D tensors;
    to apply pooling on 3D tensors,
    simply apply pooling on each channel separately.
\end{remark}

\textbf{Definition.}
Given an $H \times W \times C$ image,
we may \textbf{convolve} it using a
$Y \times X \times C$ filter (with zero-padding).
By sliding the filter along the width and height of the image,
the resulting output is an $H \times W \times 1$ image.

\begin{center}
    \frame{\includegraphics[width=0.4\textwidth]{lecture-14/convolution}}
\end{center}

(The figure above omits zero-padding,
so its output is smaller than its input.)

\begin{remark}
    There are many possible variations
    of tensor convolution that one might imagine. 
    For example, when sliding the filter,
    you might take a \emph{stride} of length (say) $2$.
    Then the output image instead has dimensions
    $\frac{H}{2} \times \frac{W}{2} \times 1$.
    \begin{center}
        \frame{\includegraphics[width=0.6\textwidth]{lecture-14/stride}}
    \end{center}
    Ordinary convolution, by default,
    uses strides of length $1$.
\end{remark}

If we want the output of convolution to have
more than one channel,
we can simply \textbf{stack} convolutions together.

\textbf{Definition.}
Given an $H \times W \times C$ image,
we may perform \textbf{stacked convolution}
using a $Y \times X \times C \times C'$ tensor of parameters.
This tensor contains $C'$ different $Y \times X \times C$ filters
stacked together.  Applying each of the $C'$ filters
to the input $H \times W \times C$ image
yields $C'$ different $H \times W \times 1$ results.
Stacking all $C'$ results together
yields an $H \times W \times C'$ output image.

\subsection{Convolutional Neural Networks}

We can now put stacked convolution and pooling together
to make a CNN.  Here's an example.

\begin{center}
    \frame{\includegraphics[width=0.9\textwidth]{lecture-14/CNN-example}}
\end{center}

Let's walk through how this CNN works.
\begin{itemize}
    \item The input image is a $32 \times 32 \times 3$ image
    of a happy face.
    \item To reach the $1^{\mathrm{st}}$ hidden layer,
    we apply a \textbf{stacked convolution} using a
    $3 \times 3 \times 3 \times 4$ tensor of parameters.
    \begin{itemize}
        \item The $1^{\mathrm{st}}$ hidden layer is then
        a $32 \times 32 \times 4$ image.
    \end{itemize}
    \item To reach the $2^{\mathrm{nd}}$ hidden layer,
    we apply \textbf{max pooling} to downsample the image
    to $16 \times 16 \times 4$.
    \item To reach the $3^{\mathrm{rd}}$ hidden layer,
    we apply a \textbf{stacked convolution} using a
    $3 \times 3 \times 4 \times 8$ tensor of parameters.
    \begin{itemize}
        \item The $3^{\mathrm{rd}}$ hidden layer is then
        a $16 \times 16 \times 8$ image.
    \end{itemize}
    \item To reach the $4^{\mathrm{th}}$ hidden layer,
    we apply \textbf{max pooling} to downsample the image
    to $8 \times 8 \times 8$.
    \item To reach the $5^{\mathrm{th}}$ hidden layer,
    we apply a \textbf{stacked convolution} using a
    $3 \times 3 \times 8 \times 16$ tensor of parameters.
    \begin{itemize}
        \item The $5^{\mathrm{th}}$ hidden layer is then
        an $8 \times 8 \times 16$ image.
    \end{itemize}
    \item To reach the $6^{\mathrm{th}}$ hidden layer,
    we apply \textbf{average pooling} to downsample the image
    to $1 \times 1 \times 16$.
    \begin{itemize}
        \item The effect of average pooling here
        is to replace every $8 \times 8$ channel
        with the average of all $64$ of its pixel values.
    \end{itemize}
    \item Finally, we interpret the $6^{\mathrm{th}}$ hidden layer
    simply as $16$ nodes, and finish with
    an ordinary \textbf{dense NN}.
    \begin{itemize}
        \item Since the output layer has $3$ nodes,
        this is a $K_{16, 3}$ bipartite graph
        with $48$ edges (hence $48$ parameters).
    \end{itemize}
\end{itemize}

And that's how a CNN works!

\subsection{Commentary on CNNs}

Some miscellaneous additional notes on CNNs:
\begin{itemize}
    \item Each stacked convolution only affects the number of channels,
    whereas each pooling only affects the width and height.
    \item It was left unstated earlier, but
    after each convolution, a bias should be added to each cell,
    then a ReLU (or similar activation).
    \begin{itemize}
        \item Each channel gets a different bias that is
        added to every cell in that channel.
        \item So the $1^{\mathrm{st}}$ stacked convolution with a
        $3 \times 3 \times 3 \times 4$ tensor of parameters
        also requires $4$ bias parameters.
    \end{itemize}
    \item Compare the CNN above with an ordinary dense NN.
    \begin{itemize}
        \item In a dense NN, reaching the first layer
        would require $(32 \times 32 \times 3) \times (32 \times 32 \times 4)$
        parameters.
        \item In this CNN, reaching the first layer
        only requires $3 \times 3 \times 3 \times 4 = 108$ parameters.
        \item Max pooling isn't implementable in an ordinary NN.
        And it requires zero parameters in this CNN.
        \item \dots something something, sparse matrices,
        something something \dots
    \end{itemize}
    \item Backpropagation is performed in the same way in CNNs
    as with NNs.
    \begin{itemize}
        \item The main difference is
        that each parameter now affects multiple pre-activations,
        not just one.
        But this just amounts to
        having to do more additions in the backwards pass.
    \end{itemize}
    \item In practice, here's what
    the learned filters for the first hidden layer might look like.
    \begin{center}
        \frame{\includegraphics[width=0.2\textwidth]{lecture-14/learned-filters}}
    \end{center}
    A lot of these filters should make sense from a
    \href{/6.300/lectures/17}{signal} \href{/6.300/lectures/18}{processing}
    perspective.
\end{itemize}

There are a lot more things that can be said
about CNN architecture, including
the endless possibilities for variations
on the overall structure.
Here's the structure of
\href{https://en.wikipedia.org/wiki/AlexNet}{AlexNet}---a
highly-successful image classifier released in 2012---for example.

\begin{center}
    \frame{\includegraphics[width=0.8\textwidth]{lecture-14/alex-net}}
\end{center}

But we won't get into too much detail on that here.

\end{document}